\documentclass[11pt,a4paper]{article}

\usepackage[margin=2.6cm]{geometry}
\usepackage{amsmath,amssymb,amsthm}
\usepackage{graphicx}
\usepackage[colorlinks=true,linkcolor=blue,citecolor=blue,urlcolor=blue]{hyperref}
\usepackage{microtype}
\usepackage{enumitem}

\theoremstyle{plain}
\newtheorem{theorem}{Theorem}[section]
\newtheorem{corollary}[theorem]{Corollary}
\newtheorem{lemma}[theorem]{Lemma}
\theoremstyle{definition}
\newtheorem{definition}[theorem]{Definition}
\newtheorem{remark}[theorem]{Remark}
\newcommand{\sgn}{\sigma}
\newcommand{\norm}[1]{\left\lVert #1\right\rVert}
\newcommand{\lamii}{\lambda_{2}}
\newcommand{\rhoD}{\rho_{d}}
\newcommand{\bdG}{A(G)}
\newcommand{\bdAs}{A_{\sgn}(G)}
\newcommand{\bbE}{\mathbb{E}}
\newcommand{\bbR}{\mathbb{R}}
\newcommand{\bbQ}{\mathbb{Q}}
\newcommand{\ind}{\mathbf 1}

\title{\bf Ramanujan signing of regular graphs:\\
a structural obstruction for cubic graphs\\[2mm]
\large and the precise remaining gap}
\author{}
\date{}

\begin{document}
\maketitle

\begin{abstract}
A \emph{signing} of a graph $G$ is a map $\sgn\colon E(G)\to\{\pm1\}$; write
$A_{\sgn}(G)$ for the resulting signed adjacency matrix and
$\beta(G)=\min_{\sgn}\norm{A_{\sgn}(G)}_{\mathrm{op}}$; put
$\rho_d=2\sqrt{d-1}$. Bilu and Linial conjectured in 2006 that every finite
simple $d$-regular graph satisfies $\beta(G)\le\rho_d$. That conjecture is
false: Xu (2026) constructed a cubic graph with $\beta(G)>2\sqrt2$, and left
open the natural refinement: does every connected regular \emph{Ramanujan}
graph admit a signing whose whole spectrum lies in
$[-\rho_d,\rho_d]$? We attack the cubic case. We prove a two-arm bottleneck
bound for cubic graphs at a cut vertex, use it to obtain
$\lambda_2\ge 20/7>2\sqrt2$ whenever two of the three components of $G-z$ have
order at least $7$, and settle the remaining family of arm orders $(5,5,\ell)$
by an exact $4\times4$ pencil analysis on a four-dimensional test space.
The outcome is:

\medskip\noindent
\textbf{Theorem A.} Every connected cubic graph with a cut vertex satisfies
$\lambda_2(\bdG)>2\sqrt2$. In particular every connected Ramanujan cubic graph is
2-connected.

We also record, and correct, two errors that arose during the investigation: that
vertex sums are invariant under vertex switching, and that
$\max_v|\sgn(v)|$ bounds $\norm{A_{\sgn}(G)}_{\mathrm{op}}$. Neither is true, the
second is refuted by $K_4$. The conjectures themselves remain open.
\end{abstract}

\section{Introduction}

Let $G$ be a finite simple undirected graph, $A(G)$ its $0/1$ adjacency
matrix, and $\mathbf 1$ the all-ones vector. A \emph{signing} is a map
$\sgn\colon E(G)\to\{\pm1\}$; the \emph{signed adjacency matrix} $A_{\sgn}(G)$
has entries $\pm1$ on $E(G)$ and $0$ elsewhere. Set
\[
  \beta(G)=\min_{\sgn}\norm{A_{\sgn}(G)}_{\mathrm{op}},\qquad
  \rho_d=2\sqrt{d-1}\quad (d\text{-regular } G).
\]
Vertex switching $\sgn\mapsto(\delta_u\sgn(\{u,v\})\delta_v)$ gives
$A_{\delta\sgn}=DA_{\sgn}D$, so the spectrum of $A_{\sgn}$ is a switching
invariant; the orbit of a signing has size $2^{n-1}$, and a set of
representatives for a connected graph on $n$ vertices with $m$ edges is obtained
by switching along a spanning tree, giving exactly $2^{m-n+1}$ classes. In
particular $\beta$ is exactly computable for cubic graphs on $n\lesssim20$
vertices.

Bilu and Linial \cite{bl06} asked whether $\beta(G)\le\rho_d$ always holds.
Marcus, Spielman and Srivastava \cite{mss15} proved the \emph{one-sided} bound
$\min_{\sgn}\lambda_{\max}(A_{\sgn})\le\rho_d$ for every $d$-regular graph,
which settles Bilu--Linial only for bipartite $G$ (where the signed spectrum is
symmetric). Recently Xu \cite{xu26} constructed a cubic graph for which every
signing has an eigenvalue outside $[-2\sqrt2,2\sqrt2]$, disproving the original
conjecture, and stated as an open question:

\begin{definition}[Ramanujan signing, conjectural \cite{xu26}]
A $d$-regular graph $G$ is \emph{Ramanujan} if every eigenvalue of $\bdG$ other
than the trivial ones lies in $[-\rho_d,\rho_d]$. \textbf{Conjecture R}: if $G$
is a connected Ramanujan graph then $\beta(G)\le\rho_d$.
\end{definition}

We consider the cubic case,
\begin{quote}\textbf{Conjecture R3.} If $G$ is connected and cubic with
$\lambda_2(\bdG)\le2\sqrt2$ then $\beta(G)\le2\sqrt2$.\end{quote}

Xu's counterexample is assembled around a cut vertex: its core graph $J$ has a
central vertex whose removal leaves three very large components. This suggests
asking whether a cut-vertex bottleneck can obstruct Ramanujan signing. The
answer is no, in a provable way.

\section{Setup at a cut vertex}

Throughout this section $G$ is connected and cubic and $z\in V(G)$ is a cut
vertex; write
\[
  G-z=C_1\;\dot\cup\;C_2\;\dot\cup\;C_3,\qquad n_i=|C_i|,\quad n_1\le n_2\le n_3,
\]
and let $t_i$ be the unique vertex of $C_i$ adjacent to $z$.

\begin{lemma}
$G-z$ has exactly three components, each of odd order at least $5$, each joined
to $z$ by exactly one edge. Moreover $2|E(C_i)|=3n_i-1$.
\end{lemma}

\begin{proof}
Every component of $G-z$ contains a neighbour of $z$, since otherwise $G$ is
disconnected. As $\deg(z)=3$ and the components partition the neighbours of $z$,
there are exactly three components and each is attached to $z$ by one edge.
Summing degrees over $C_i$ gives $3n_i=2|E(C_i)|+1$, so $n_i$ is odd and
$|E(C_i)|=(3n_i-1)/2$. As $C_i$ is simple,
$\tfrac{3n_i-1}{2}\le\tfrac{n_i(n_i-1)}{2}$, i.e.\ $n_i^2-4n_i+1\ge0$, whence
$n_i\ge4$ (the roots are $2\pm\sqrt3$); oddness gives $n_i\ge5$.
\end{proof}

Two families of vectors orthogonal to $\mathbf 1$ will be used:
\[
  p_i=\ind_{C_i}-n_i e_z,\qquad
  q_i=e_{t_i}-\tfrac1{n_i}\ind_{C_i}.
\tag{2.1}
\]
Indeed $\langle p_i,\mathbf 1\rangle=n_i-n_i=0$ and
$\langle q_i,\mathbf 1\rangle=1-\frac1{n_i}n_i=0$.

\section{The bottleneck bound}

Since $G$ is connected, $\mathbf 1$ spans the top eigenspace of $\bdG$ and
\[
  \lambda_2(\bdG)=\max\Big\{\tfrac{x^{\mathsf T}\bdG x}{x^{\mathsf T}x}:
  0\neq x\perp\mathbf 1\Big\}.
\tag{3.1}
\]

\begin{theorem}[two-arm bottleneck bound]
$\lambda_2(\bdG)\ge\max_{i<j}\left(3-\dfrac{n_i^2+n_j^2}{n_in_j(n_i+n_j)}\right)$,
and in particular $\lambda_2(\bdG)\ge\max_{i<j}\left(3-\dfrac1{\max(n_i,n_j)}\right)$.
\end{theorem}

\begin{proof}
Fix two components with orders $a=n_i$, $b=n_j$ and set
$x=\tfrac1a\ind_{C_i}-\tfrac1b\ind_{C_j}$. Then $\langle x,\mathbf 1\rangle=0$
and, because $C_i,C_j$ are disjoint and $x$ vanishes elsewhere,
\[
  \norm{x}^{2}=\tfrac1a+\tfrac1b.
\]
There are no edges between $C_i$ and $C_j$, and all edges leaving them go to $z$
where $x_z=0$; using $2|E(C_k)|=3n_k-1$,
\[
  x^{\mathsf T}\bdG x=\frac{3a-1}{a^2}+\frac{3b-1}{b^2}.
\]
Putting numerator over $a^2b^2$ and denominator over $ab$,
\[
  \frac{x^{\mathsf T}\bdG x}{\norm{x}^{2}}
  =\frac{b^2(3a-1)+a^2(3b-1)}{ab(a+b)}
  =\frac{3ab^{2}-b^{2}+3a^{2}b-a^{2}}{ab(a+b)}
  =3-\frac{a^{2}+b^{2}}{ab(a+b)},
\]
the last step being the identity
$3ab^{2}+3a^{2}b-a^{2}-b^{2}=3ab(a+b)-(a^{2}+b^{2})$. By (3.1) this
quotient is at most $\lambda_2(\bdG)$; maximising over the three pairs gives the
claim. Finally, writing $t=a/b\in(0,1]$,
$\frac{a^2+b^2}{ab(a+b)}=\frac1a\cdot\frac{t+1/t}{1+t}\le\frac1a$ because
$\frac{t+1/t}{1+t}\le1\iff\frac1t\le1$.
\end{proof}

\begin{corollary}
If at least two of $n_1,n_2,n_3$ are at least $7$, then
$\lambda_2(\bdG)\ge\frac{20}{7}>2\sqrt2$.
\end{corollary}

\begin{proof}
Apply the theorem to those two arms; the smaller order $a\ge7$ gives
$\lambda_2\ge3-\frac1a\ge3-\frac17=\frac{20}{7}$, and
$\frac{400}{49}>8$.
\end{proof}

Thus only the configuration $n_1=n_2=5$ (so $n_3=\ell$ odd, $\ge5$) survives.

\section{The residual family $(5,5,\ell)$}

Consider the four-dimensional space
\[
  W=\operatorname{span}\{P,Q,R,S\},\quad
  P=p_1+p_2,\quad Q=q_1+q_2,\quad R=p_3,\quad S=q_3,
\]
contained in $\mathbf 1^\perp$ by (2.1). Hence
$\lambda_2(\bdG)\ge\lambda_{\max}(M_\ell,K_\ell)$ where
$M_\ell=p^{\mathsf T}\bdG p$, $K_\ell=p^{\mathsf T}p$.

We derive the pencil from three combinatorial identities. Write a vector as
value $c_i$ on $C_i\setminus\{t_i\}$, value $c_i+d_i$ at $t_i$, and value
$z_x$ at $z$. Then
\begin{enumerate}[leftmargin=2em,itemsep=2pt]
\item[(i)] On one arm,
$\sum_{u\in C_i}x(u)y(u)=(n_i-1)c^xc^y+(c^x+d^x)(c^y+d^y)$.
\item[(ii)] Arm-internal adjacency form. Since every vertex of $C_i$ except
$t_i$ has degree $3$ inside $C_i$ and $t_i$ has degree $2$,
$x^{\mathsf T}\bdG y|_{C_i}=\sum_{u\in C_i}x(u)\sum_{v\in C_i,\,v\sim u}y(v)$.
Grouping by $v\neq t_i$ gives $3c^x\sum_{v\neq t_i}y(v)$ except that for
$v\sim t_i$ the value $x(t_i)$ replaces $c^x$ once, contributing $d^xy(v)$;
the pair $\{t_i,u\}$ is then counted twice, the second time contributing
$y(t_i)\sum_{u\sim t_i,\,u\neq t_i}x(u)=2(c^y+d^y)c^x$. Hence
\[
   x^{\mathsf T}\bdG y|_{C_i}=3(n_i-1)c^xc^y+2d^xc^y+2c^x(c^y+d^y).
\]
\item[(iii)] Each of the three edges $zt_i$ contributes
$z_x\,y(t_i)+x(t_i)\,z_y$.
\end{enumerate}

Substituting $n_1=n_2=5$, $n_3=\ell$ yields
\[
 M_\ell=\begin{pmatrix}
 -12 & -\tfrac{88}{5} & -2\ell-10 & -10+\tfrac{10}{\ell}\\
 -\tfrac{88}{5} & -\tfrac{12}{25} & -\tfrac{8\ell}{5} & 0\\
 -2\ell-10 & -\tfrac{8\ell}{5} & \ell-1 & -\ell+\tfrac1\ell\\
 -10+\tfrac{10}{\ell} & 0 & -\ell+\tfrac1\ell & -\tfrac1\ell-\tfrac1{\ell^2}
 \end{pmatrix},\quad
 K_\ell=\begin{pmatrix}
 110 & 0 & 10\ell & 0\\
 0 & \tfrac85 & 0 & 0\\
 10\ell & 0 & \ell^2+\ell & 0\\
 0 & 0 & 0 & 1-\tfrac1\ell
 \end{pmatrix}.
\]
Note that the pencil depends only on the arm orders, not on the internal
structure of the arms, so the bound is uniform. This computation is
machine-checked against matrices built from explicit coordinate vectors on
concrete three-arm graphs.

Let $D_\ell=M_\ell-2\sqrt2\,K_\ell$. The principal minor of $D_\ell$ on the
pair $(P,R)$ is
\[
  \det D_\ell[\{P,R\}]
  =-4(\ell+11)\Big((14\sqrt2-19)\ell-5\sqrt2+2\Big).
\tag{5.1}
\]
Since $14\sqrt2-19>0$ and
$\frac{5\sqrt2-2}{14\sqrt2-19}=\frac{67\sqrt2+102}{31}\approx6.347$, the minor
is negative for every integer $\ell\ge7$. A symmetric $2\times2$ matrix with
negative determinant is indefinite, so $D_\ell$ admits a direction $a$ with
$a^{\mathsf T}D_\ell a>0$, i.e.\ $a^{\mathsf T}M_\ell a>2\sqrt2\,a^{\mathsf T}K_\ell a$,
and $a^{\mathsf T}K_\ell a>0$ because $K_\ell$ is a Gram matrix of linearly
independent vectors. Hence $\lambda_{\max}(M_\ell,K_\ell)>2\sqrt2$ for all
$\ell\ge7$.

For $\ell=5$ take the rational vector
$a=(-\tfrac35,\tfrac15,\tfrac65,-\tfrac25)^{\mathsf T}$ in the basis
$(P,Q,R,S)$. Direct substitution gives
\[
  a^{\mathsf T}M_5a-2\sqrt2\,a^{\mathsf T}K_5a
  =\frac{19584-13740\sqrt2}{625}>0,
\]
because $19584^2-2\cdot13740^2=6\,017\,856>0$ and $19584>0$.

\section{Main theorem}

\begin{theorem}
Every connected cubic graph with a cut vertex satisfies
$\lambda_2(\bdG)>2\sqrt2$.
\end{theorem}

\begin{proof}
If $n_2\ge7$, Corollary 1 applies. Otherwise $n_1=n_2=5$ and $n_3=\ell\ge5$ is
odd, and Section 4 applies ($\ell\ge7$ or $\ell=5$).
\end{proof}

\begin{corollary}
Every connected Ramanujan cubic graph is 2-connected.
\end{proof}
By Theorem A, a connected cubic graph with a cut vertex has a non-trivial
eigenvalue of modulus $>2\sqrt2=\rho_3$, so it is not Ramanujan. For a
connected graph the top eigenspace of $\bdG$ is $\operatorname{span}\{\mathbf 1\}$,
so the second largest eigenvalue is non-trivial.
\end{proof}

\section{Computational evidence}

Tables produced by \texttt{experiments/run\_all\_experiments.py}; $\beta(G)$ is
exact (all $2^{n/2+1}$ switching classes are enumerated).

\begin{center}
\small
\begin{tabular}{lrrrr}
\toprule
$G$ & $n$ & $\lambda_2(\bdG)$ & Ramanujan & $\beta(G)$\\
\midrule
$K_4$ & 4 & $-1.000000$ & yes & $2.236068$\\
$K_{3,3}$ & 6 & $0.000000$ & yes & $2.000000$\\
triangular prism $C_3\times K_2$ & 6 & $1.000000$ & yes & $2.236068$\\
cube $C_4\times K_2$ & 8 & $1.000000$ & yes & $1.732051$\\
pentagonal prism & 10 & $1.618034$ & yes & $2.236068$\\
hexagonal prism & 12 & $2.000000$ & yes & $2.000000$\\
Heawood & 14 & $1.414214$ & yes & $2.449490$\\
Möbius--Kantor $GP(8,3)$ & 16 & $1.732051$ & yes & $2.414214$\\
\bottomrule
\end{tabular}
\end{center}

All eight satisfy $\beta\le2\sqrt2$, as Conjecture R3 predicts. Eight graphs
are evidence, not a proof.

\section{Discussion, limitations, open questions}

\subsection*{Two claims made and withdrawn during this work}

\emph{(i) Vertex sums are not switching-invariant.} It is true that
$A_{\delta\sgn}=DA_{\sgn}D$, hence the spectrum is preserved, but
$A_{\sgn}\mathbf 1=(v\mapsto\sgn(v))$ gives
$A_{\delta\sgn}\mathbf 1=DA_{\sgn}(\delta\mathbf 1)$, whose $v$-th entry is
$\delta(v)\sum_{w\sim v}\sgn(\{v,w\})\delta(w)$, generally not $\pm\sgn(v)$.
On $C_4$ with $\sgn\equiv+1$ (all vertex sums $+1$), switching by
$\delta=(1,1,-1,-1)$ gives all vertex sums $0$.

\emph{(ii) $\max_v|\sgn(v)|$ does not bound $\norm{A_{\sgn}}_{\mathrm{op}}$.}
For $K_4$ with $\sgn(\{1,2\})=\sgn(\{3,4\})=-1$ all vertex sums are $\pm1$
while the eigenvalues of $A_{\sgn}$ are $3,-3,1,1$, so the norm is $3$. The
true optimum is $\beta(K_4)=\sqrt5$, attained by signings with larger vertex
sums. Consequently signings coming from perfect matchings --- the natural first
idea --- give no bound at all.

Both errors were caught by an explicit cross-check against concrete graphs and
are documented in the repository's research log.

\subsection*{Limitations}

Theorem A rules out cut vertices but says nothing about 2-vertex cuts, small
bridges, or finite bottlenecks. The case $d\ge4$ is untouched. The bottleneck
estimate itself is an elementary instance of the classical Cheeger-type
argument; what is proved here is the precise statement and its use in the
residual case. We do not claim historical priority for Theorem A.

\subsection*{Open questions}

\begin{enumerate}[leftmargin=2em]
\item Do Ramanujan cubic graphs admit signings with
$\norm{A_{\sgn}}\le2\sqrt2$ (Conjecture R3)?
\item Is every $3$-connected cubic graph $2$-realisable, i.e.\ $\beta(G)\le2\sqrt2$?
\item What is the smallest cubic graph with $\beta(G)>2\sqrt2$? Xu's example has
astronomically many vertices; exhaustive search is cheap with the switching-class
machinery (there are only $4060$ connected cubic graphs on $16$ vertices and
$2^9$ switching classes each).
\item Is the bound $3-\frac{a^2+b^2}{ab(a+b)}$ of Theorem 1 attained, and for
which $(a,b)$?
\end{enumerate}

\begin{thebibliography}{9}
\bibitem{bl06} Y.~Bilu and N.~Linial. Lifts, discrepancy and nearly optimal
spectral gap. \emph{Combinatorica} 26 (2006) 495--519.

\bibitem{xu26} Z.~Xu. A 3-regular counterexample to the Bilu--Linial signing
conjecture. arXiv:2609.15591v4 (2026).

\bibitem{mss15} A.~W.~Marcus, D.~A.~Spielman and N.~Srivastava. Interlacing
families I: bipartite Ramanujan graphs of all degrees. \emph{Annals of
Mathematics} 182 (2015) 307--325.

\bibitem{rs21} N.~Ravichandran and N.~Srivastava. Asymptotically optimal
multi-paving. \emph{IMRN} 2021, 10908--10940.

\bibitem{hps18} R.~Hall, D.~Puder and M.~Sawin. Ramanujan coverings of graphs.
\emph{Advances in Mathematics} 323 (2018) 367--397.

\bibitem{lz26} F.~Lin and H.~Zhou. Improved bounds for the Bilu--Linial
conjecture via spectral recovery from mixed determinantal polynomials.
arXiv:2609.15715 (2026).
\end{thebibliography}

\end{document}